%%%%%%%%%%  scan idx 177  |  folio 164  %%%%%%%%%%
\[
H^1(S,F)=0\quad
\text{pour tout }F\in\Ob\mathcal C'_S.
\]
Alors :
\begin{enumerate}[label=\alph*)]
\item Toute suite exacte de $\mathcal C'_S$ :
\[
0\longrightarrow E\longrightarrow F\longrightarrow G\longrightarrow0,
\]
telle que $\underline{\Hom}(G,E)
\in\Ob\mathcal C'_S$ et que $G$ soit localement facteur direct d'un Module libre de type fini, est scindée. La catégorie $\mathcal C_{0S}$ est relevable dans $\mathcal C'_S$ (I 1.2).
\item Notons $\Mod(A(S))$ la catégorie des $A(S)$-modules. Le foncteur
\[
\Gamma_S:\mathcal C'_S\longrightarrow\Mod(A(S))
\]
est exact. En outre :
\begin{enumerate}[label=(\roman*)]
\item Si tout objet $F$ de $\mathcal C'_S$ admet une présentation finie globale
\[
A_S^p\longrightarrow A_S^q\longrightarrow F\longrightarrow0,
\]
le foncteur $\Gamma_S:\mathcal C'_S\longrightarrow\Mod(A(S))$ est pleinement fidèle et a pour image essentielle la sous-catégorie pleine de $\Mod(A(S))$ formée des $A(S)$-modules de présentation finie ;
\item Si $\mathcal C'_S$ contient les $A_S$-Modules libres, si tout objet $F$ de $\mathcal C'_S$ admet une présentation globale
\[
A_S^{(I)}\longrightarrow A_S^{(J)}\longrightarrow F\longrightarrow0,
\]
et si le foncteur $\Gamma_S$ commute aux limites inductives filtrantes\footnote{hypothèse vérifiée si $S$ est localement de type fini (SGA 4 VI 1.2) ou si $S$ est le topos des faisceaux d'ensembles sur un espace compact et $S$ l'objet final (T.F II 3.10.1).}, le foncteur $\Gamma_S:\mathcal C'_S\longrightarrow\Mod(A(S))$ est une équivalence de catégories.
\end{enumerate}
\end{enumerate}

\emph{Preuve.} a) Comme $G$ est localement facteur direct d'un libre de type fini, la suite spectrale
\[
E_2^{pq}=H^p(S,\underline{\Ext}^q(G,E))\Longrightarrow\Ext^{*}(G,E)
\]
dégénère (I 1.3.1) et donne un isomorphisme
\[
H^1(S,\underline{\Hom}(G,E))\simeq\Ext^1(G,E),
\]
